% Presentación - Estado de Avance (10 minutos).
% Rúbrica: Representación y restricciones (20 %), Algoritmo (60 %),
% Conclusiones (10 %), Otros (10 %: referencias, gráficos, duración).
\documentclass[aspectratio=169, 10pt]{beamer}
\usepackage[utf8]{inputenc}
\usepackage[T1]{fontenc}
\usepackage[spanish]{babel}
\usepackage{amsmath}
\usepackage{graphicx}
\usepackage{xcolor}
\usepackage[cache=true]{minted}
\setminted{fontsize=\normalsize, style=vs, bgcolor=gray!8, linenos=false, breaklines=true, highlightcolor=yellow!45}

\usetheme{Madrid}
\setbeamertemplate{bibliography item}[text]

\title{Segmentación multiumbral de imágenes MRI}
\subtitle{Greedy + Tabú Search --- Estado de avance}
\author{Matías Peñaloza}
\institute{Inteligencia Artificial 2026-2}
\date{\today}

\begin{document}

\begin{frame}
  \titlepage
\end{frame}

% ---------------------------------------------------------------- 0:00-1:30
\begin{frame}{Problema, representación y restricciones}
  \textbf{Problema:} Dado un histograma de 256 niveles y un $k$ fijo, elegir $k$
  umbrales que maximicen la varianza entre clases segun el criterio de
  Otsu~\cite{Otsu1979}.

  \textbf{Representación de una solución:}
  $t = (t_1, t_2, \ldots, t_k)$ un vector de $k$ umbrales

  \textbf{Dominios:} $\forall\, t_i \in \{1, \ldots, 254\} \subset \mathbb{Z}$

  \textbf{Restricciones:}
  \begin{itemize}
    \item Cota inferior: $1 \leq t_1$
    \item Monotonia estricta: $t_i < t_{i+1}\ \forall\, i = 1, \ldots, k-1$
    \item Cota superior: $t_k \leq 254$
  \end{itemize}

\end{frame}

\begin{frame}{Problema, representación y restricciones}
  \textbf{Implementación:}
  \begin{itemize}
    \item Instancia:
    \begin{itemize}
      \item \texttt{std::array<unsigned long long, 256> histograma}
      \item \texttt{int k}
    \end{itemize}
    \item Solución: \texttt{std::vector<int> umbrales}
    \begin{itemize}
      \item No garantiza orden en los umbrales (\texttt{std::set}) pero es mejor en la reparación
      \item No tiene tamaño fijo a pesar de un $k$ fijo (\texttt{std::array}) pero es mejor en la generación
      \item Garantiza umbrales enteros
    \end{itemize}
  \end{itemize}
\end{frame}

\begin{frame}[fragile]{Problema, representación y restricciones}
  \textbf{Manejo de restricciones --- Generación de solución inicial:}

  Greedy busca secuencialmente los umbrales, se inserta uno artificial 0 al inicio del vector
\begin{minted}[highlightlines={2,5}]{cpp}
std::vector<int> generarSolucionInicial(...) {
  std::vector<int> umbrales = {0};
  ..
  greedy(umbrales, f, instancia.k, tabla, gen);
  umbrales.erase(umbrales.begin());
  return umbrales;
}
\end{minted}
\end{frame}

\begin{frame}[fragile]{Problema, representación y restricciones}
  \textbf{Manejo de restricciones --- Generación de solución inicial:}

  En cada iteración se setea la cota inferior con respecto al umbral anterior que conserva la monotonía y una cota superior hasta el máximo posible para cada umbral.
\begin{minted}[highlightlines={4,6}]{cpp}
void greedy(std::vector<int> &umbrales, double &facc, int k, const otsu::TablaOtsu &tabla, std::mt19937 &gen) {
  int size = static_cast<int>(umbrales.size());
  ..
  int cotainf = umbrales.back() + 1, cota_sup = 254 - k + size;
  ..
  for (int um = cotainf; um <= cota_sup; um++) {
    ..
  };
  ..
  greedy(umbrales, facc, k, tabla, gen);
}
\end{minted}
\end{frame}

\begin{frame}[fragile]{Problema, representación y restricciones}
  \textbf{Manejo de restricciones --- Reparación de la solución:}

  Tabu search filtra los valores posibles para los vecinos entre la cota superior e inferior. Asegura monotonía estricta usando pasos de $\pm 1$ y filtrando vecinos que ya existan en el vector de umbrales ($t_i + \text{mov} \notin t$).
\begin{minted}[fontsize=\scriptsize, highlightlines={7-10,15-16}]{cpp}
std::vector<std::tuple<int, int>> generarVecinos(std::vector<int> &actual, int size) {
  std::vector<std::tuple<int, int>> vecinos;

  auto agregarVecinos = [&actual, &vecinos, &size](int mov) {
    for (int i = 0; i < size; i++) {
      int val = actual[i] + mov;
      if (val < 1 || val > 254)
        continue;
      if (std::find(actual.begin(), actual.end(), val) != actual.end())
        continue;
      ..
    }
  };

  agregarVecinos(1);
  agregarVecinos(-1);
  return vecinos;
}
\end{minted}
\end{frame}

% ---------------------------------------------------------------- 1:30-3:00
\begin{frame}{Lectura, función de evaluación y salida}
  % TODO: lectura/validación de instancias, tablas de sumas acumuladas
  %       precalculadas (O(256) una vez, luego O(k) por evaluación)
  %       \cite{Liao2001Fast}, archivos de salida; evidencia de que funciona
  %       en 30/30.
\end{frame}

% ---------------------------------------------------------------- 3:00-5:00
\begin{frame}{Solución inicial: Greedy}
  % TODO: criterio final y por qué falló la primera versión (historia del
  %       diseño); gráfico de brecha por criterio.
\end{frame}

% ---------------------------------------------------------------- 5:00-7:30
\begin{frame}{Tabú Search: movimientos y diseño}
  % TODO: vecindario ±1, lista tabú y tenencia, aspiración, caso "todos
  %       tabú"; justificación de cada decisión~\cite{Hammouche2010Comparative}.
\end{frame}

% ---------------------------------------------------------------- 7:30-9:00
\begin{frame}{Calidad y tiempos}
  % TODO: tabla/gráfico de brecha contra el óptimo exacto (k=2, 4 y
  %       k=6 en una imagen) y tiempos en escala logarítmica.
\end{frame}

% ---------------------------------------------------------------- 9:00-10:00
\begin{frame}{Conclusiones y trabajo restante}
  % TODO: conclusiones acordes a lo expuesto y trabajo que falta.
\end{frame}

\begin{frame}[allowframebreaks]{Referencias}
  \nocite{Otsu1979, Hammouche2010Comparative, Liao2001Fast}
  \bibliographystyle{plain}
  \bibliography{Referencias}
\end{frame}

\end{document}
